\section{Introduction}\label{sec:introduction}

When a prediction market prices a candidate at 80 cents, the price is read as an 80\% chance of winning. Since a US
court allowed Kalshi to list election contracts in 2024, such prices have moved onto regulated exchanges, and more than
\$2 billion traded on markets about the 2024 presidential race in its last five weeks alone
\citep{clinton2025}.\footnote{In September 2024 the US District Court for the District of Columbia ruled for Kalshi in
\emph{KalshiEX LLC v.\ Commodity Futures Trading Commission}. The regulator withdrew its appeal in May 2025.} Whether
these prices can be read as probabilities decides how much weight the forecasts built on them deserve.

Betting markets suggest they cannot, quite. At racetracks, bettors overpay for long shots and underpay for favourites
\citep{griffith1949,thaler1988}, most likely because they misjudge small probabilities \citep{snowberg2010}. A market
with this favourite--longshot bias is \emph{underconfident}: its prices sit too close to even odds, so favourites win
more often than their price says. Prediction markets show the same tendency. Their prices are well calibrated shortly
before a market closes but biased for events further away \citep{page2013}. On Kalshi, expensive contracts earn small
positive returns while cheap ones lose \citep{burgi2025}, and across about 327{,}000 contracts on Kalshi and
Polymarket, political markets are underconfident at nearly every time horizon \citep{le2026}. That political prediction
markets are slightly underconfident is close to settled.

Two features of this evidence leave a question open. First, it pools all political markets, although elections differ
in exactly the ways that should matter. A general election has party nominees, regular polls, heavy coverage and party
loyalties that make most results predictable. A primary is fought within one party, often between several candidates,
with few polls and thin trading, and party labels tell traders nothing. There is no reason to expect the same bias in
both. Second, the evidence is retrospective. Patterns are found and reported in the same data, after many reasonable
choices about which markets to group and which to drop, and may not survive the next election \citep{gelman2014}.
Pre-registration fixes the analysis, and what each possible result would mean, before the data that decide it exist
\citep{nosek2018,olken2015}. With markets this is rarely possible, but elections, whose dates are known in advance,
allow it.

We study every two-candidate US contest settled on Kalshi or Polymarket between November 2023 and September 2026, 152
on Kalshi and 64 on Polymarket, each priced once at a fixed time on election day before the polls close. One number,
the \emph{stretch}, measures whether favourites won more often (above 1) or less often (below 1) than their prices
implied. General elections show the familiar underconfidence: on 5 November 2024 the stretch was 1.97 on Kalshi and
1.83 on Polymarket. Kalshi's primaries go the other way (0.32), while Polymarket's were close to calibrated (0.88).
Pooled, Kalshi's markets even look overconfident (0.58), which shows how much a single figure for politics can hide.

Most general-election contests in these data come from 2024 and most primaries from 2026, so the split has two possible
explanations: primaries are traded differently, or the underconfidence of 2024 faded over time. The US midterms of
3 November 2026, general elections held in 2026, separate the two. Before the vote we froze and publicly archived the
code, the data, the rules for selecting and pricing contests, the test, and the reading attached to each possible
outcome \citep{kirillov2026code,kirillov2026data}.

This preliminary version reports the test on the contests settled by 10 November 2026: 187 on Kalshi and 145 on
Polymarket, with 40 still open. Polymarket's general-election markets remained underconfident (stretch 1.66, 95\%
interval 1.21 to 2.91) and consistent with 2024, the pre-registered outcome that supports primaries being traded
differently. Kalshi's result is inconclusive (1.32, 0.88 to 2.35). Exploratory analyses suggest why: on the 141 races
both platforms listed, they priced almost identically and were equally underconfident, and Kalshi's weaker figure comes
from lower-profile races that only it listed. Buying every midterm favourite, after costs, earned small returns
that cannot be told apart from zero.

Section~\ref{sec:related} reviews related work, Sections~\ref{sec:data} and~\ref{sec:method} describe the data and the
pre-registered test, Section~\ref{sec:results} reports the results and Section~\ref{sec:backtest} the backtest.
Sections~\ref{sec:limitations} to~\ref{sec:conclusion} cover limitations, discussion and conclusions.
